This project investigates whether sparse autoencoders (SAEs) expose
interpretable internal features associated with feedback uptake in a
controlled Solver--Critic pipeline. A Gemma 3 4B instruction-tuned Solver
answers a mixture of MuSiQue and synthetic multi-hop questions. A Qwen3 4B
instruction-tuned Critic first forms an answer without seeing the Solver's
answer; the Solver then receives both candidate answers. The primary analysis
focuses on episodes labeled Adoption or Retention according to whether the
Solver adopts the Critic's answer or retains its own. The final core corpus
contains 91,063 episodes across the training, validation, and test partitions,
and a separately
collected intervention partition contains 10,105 episodes. We cache the
Solver's Layer 25 residual-stream activation at the final prompt token,
represent it with a custom $L_1$ SAE and pretrained Gemma Scope SAE variants,
and fit regularized probes on training data to identify sparse features that
distinguish Adoption from Retention on validation data. Feature evidence
includes activating examples, held-out predictive performance, permutation
tests, matched random-feature controls, and interventions on a partition sealed
during feature selection. In an exploratory experiment on a nearly balanced
995-episode intervention cohort (498 Adoption and 497 Retention episodes),
targeted edits using an LLM-edited pretrained SAE produced pooled flip rates of
11.56--13.37\%, compared with 2.78--3.79\% for matched random-feature edits.
All three targeted-versus-random differences survived Holm correction. These
results show that the selected directions can change the judged feedback-uptake
outcome more often than matched random directions under this intervention
protocol. Because the LLM-edited SAE experiment was exploratory and its edits
also increased off-target outputs, the result does not by itself establish a
simple human-readable meaning for each selected feature. The central
contribution is an evidence-based framework for evaluating feature
interpretability by combining semantic inspection, held-out prediction, and
controlled causal testing.
